\subsection{Elección de la temporada}
La selección de las temporadas responde a las necesidades de cada fase del proyecto. En la primera fase se busca ampliar el conocimiento sobre el desempeño de cada jugadora durante los intervalos en que comparte cancha con una oponente determinada. Incorporar temporadas anteriores permite recuperar antecedentes de estas coincidencias y puede aumentar las observaciones disponibles, siempre que ambas jugadoras hayan participado en esos periodos. No obstante, este beneficio depende del historial de cada pareja y requiere considerar los cambios en su desempeño y contexto.

En la segunda fase, el interés se centra en el rendimiento colectivo de los quintetos de la temporada objetivo. La información histórica debe utilizarse con mayor cautela, pues las plantillas no permanecen necesariamente iguales entre temporadas. Como referencia, a partir de la tabla oficial de comparación de plantillas publicada por Minnesota Lynx el 14 de junio de 2025, se calculó que los 12 equipos con participación en 2024 conservaron, en promedio, 5.33 jugadoras de su plantilla final de ese año. [X] Esta variación aconseja no considerar automáticamente equivalentes las configuraciones colectivas de distintas temporadas.

Por ello, se propone utilizar varias temporadas para enriquecer el análisis individual de la primera fase y concentrar el análisis de quintetos en la temporada objetivo. Esta decisión se complementa con un cuadro comparativo de las ventajas y desventajas de ambas alternativas.

\begin{longtable}{@{}p{2.15cm}p{5.35cm}p{5.35cm}@{}}
\caption{Instrumento de ayuda a la elección de datos en determinada temporada por fase.}\\
\toprule
\textbf{Fase}&\textbf{Una temporada}&\textbf{Más de una temporada}\\\midrule\endfirsthead
\toprule\textbf{Fase}&\textbf{Una temporada}&\textbf{Más de una temporada}\\\midrule\endhead
Jugadora vs. jugadora&
\textbf{Ventajas:} Mayor cercanía al desempeño, rol y contexto competitivo de la temporada objetivo. Reduce la mezcla de etapas diferentes de una misma jugadora. Facilita la preparación y comparación de los datos.

\textbf{Desventajas:} Algunas parejas oponentes pueden acumular pocas posesiones o minutos compartidos. Esa escasez puede producir estimaciones con mayor incertidumbre y sensibilidad a partidos atípicos. Limita la cobertura de parejas y puede favorecer el sobreajuste si se emplean modelos complejos.&
\textbf{Ventajas:} Puede aumentar las posesiones, partidos y contextos observados para las parejas con historial compartido. Una mayor muestra comparable puede reducir la variabilidad de las estimaciones y permitir distinguir patrones persistentes de fluctuaciones ocasionales. Amplía el historial individual disponible.

\textbf{Desventajas:} Introduce riesgo de sesgo temporal por cambios de edad, condición física, rol, equipo, compañeros o sistema táctico. Las parejas no necesariamente acumulan más enfrentamientos por añadir temporadas. Requiere ponderación temporal, regularización y control del contexto; los registros adicionales tampoco son necesariamente independientes.\\\midrule
Mejor quinteto&
\textbf{Ventajas:} Permite estimar el desempeño colectivo dentro de una plantilla y un contexto temporal más cercanos al objetivo de recomendación. Facilita interpretar la participación conjunta de las jugadoras y su rendimiento frente a los rivales de esa temporada. Evita atribuir directamente al equipo actual resultados obtenidos con plantillas anteriores.

\textbf{Desventajas:} Los quintetos poco utilizados pueden presentar estimaciones inestables. Las alineaciones nunca observadas requieren un modelo capaz de generalizar a partir de información individual o de combinaciones parciales. Persisten cambios dentro de la temporada por lesiones, transferencias o ajustes tácticos.&
\textbf{Ventajas:} Puede aportar más observaciones de quintetos recurrentes y de parejas o tríos que permanecen juntos. Permite estudiar continuidad o cambios del rendimiento y entrenar modelos con una variedad más amplia de configuraciones.

\textbf{Desventajas:} El aumento total de registros no implica más evidencia para cada quinteto actual. Mezclar temporadas puede confundir la contribución de una combinación con cambios de plantilla, entrenador, estilo o rivales. Las alineaciones históricas pueden tener poca representatividad para la temporada objetivo; se necesita modelar explícitamente esa heterogeneidad.\\
\bottomrule
\end{longtable}

\subsection{Alternativas de ponderación de los datos}
\label{subsec:alternativas-ponderacion}
La elección de las temporadas delimita la información disponible, pero no determina por sí misma la influencia que debe tener cada antecedente en la caracterización de las jugadoras. En la fase 1 se considerarán alternativas que distingan la antigüedad de las observaciones, el macrorrol de las participantes y las características estadísticas empleadas. Su comparación permitirá establecer si una mayor diferenciación mejora la estabilidad de los perfiles y la predicción del rendimiento, sin asumir que la alternativa más compleja será superior.

Se distinguirán tres operaciones: la ponderación de observaciones históricas, la comparación condicionada por posición y la ponderación de características para el agrupamiento. Aunque pueden combinarse, no son equivalentes. Separar modelos o utilizar referencias por posición constituye una forma de contextualización; solo representa una ponderación en sentido estricto cuando se modifica explícitamente la contribución de registros o variables. Asimismo, estandarizar las escalas no equivale a aprender la importancia predictiva de cada indicador \cite{bishop2006pattern,huang2005variableweighting}.

Esta distinción resulta relevante porque las exigencias estadísticas varían entre posiciones. Gómez Ruano et al. analizaron jugadoras de la WNBA de 2005 y encontraron diferencias entre bases y pívots en asistencias, triples, rebotes y tapones, con valores intermedios para las aleros \cite{gomez2007posicioneswnba}. Este antecedente justifica examinar referencias específicas por macrorrol, pero sus resultados no se trasladarán automáticamente a las temporadas del proyecto. Un registro elevado de rebotes para un ala podrá describirse como un rasgo de su perfil estadístico; no permitirá inferir su estatura ni una asignación defensiva concreta.

La unidad de análisis continuará siendo una observación dirigida por partido, jugadora focal y jugadora rival, delimitada a los intervalos de coincidencia en cancha. Las dos perspectivas de una pareja permanecerán vinculadas al mismo partido. Los identificadores, la fecha y los nombres de equipos conservarán su función de identificación y contexto; no se multiplicarán por ponderaciones. Las métricas de entrada se construirán con antecedentes anteriores al momento de predicción. La ponderación no convertirá la coincidencia en cancha en evidencia de defensa directa.

La Tabla~\ref{tab:tipos-ponderacion-fase1} sintetiza los tres ejes que originan las alternativas posteriores. En este apartado se presentan sus fundamentos y aplicación prevista; no se establecen todavía coeficientes definitivos ni expresiones de cálculo.

\begin{table}[htbp]
\centering
\small
\caption{Ejes considerados para el tratamiento diferenciado de la información en la fase 1.}
\label{tab:tipos-ponderacion-fase1}
\begin{tabularx}{\textwidth}{@{}p{2.35cm}X@{}}
\toprule
\textbf{Eje} & \textbf{Interpretación en el proyecto}\\
\midrule
Temporada & Modificar la influencia de las observaciones según su antigüedad, conservando su exposición y contexto.\\
Posición & Establecer referencias o parámetros condicionados a Guard, Wing y Big, sin confundir macrorrol con perfil funcional.\\
Variable & Diferenciar la contribución de características al clustering. Como experimento separado, evaluar descuentos temporales distintos por métrica.\\
\bottomrule
\end{tabularx}
\end{table}

% La profundidad de numeración se amplía solo para las subdivisiones de 6.1.2.
\edef\profundidadAntesPonderacion{\number\value{secnumdepth}}
\setcounter{secnumdepth}{3}

\subsubsection{Ponderación por temporada}
\label{subsubsec:ponderacion-temporada}
La información antigua puede aportar observaciones adicionales, pero también describir una etapa distinta de la jugadora. Se considerará, por ello, una reducción progresiva de su influencia. Como antecedente metodológico, Dixon y Coles incorporaron ponderación temporal en un modelo de resultados de fútbol \cite{dixon1997scores}. Se recupera el principio de descontar información antigua, no sus parámetros ni una eficacia demostrada para la WNBA. En baloncesto femenino, Chen et al. advierten que las composiciones y características de las alineaciones pueden variar entre torneos \cite{Chen17122023}.

En el dataset, la temporada y la fecha permitirán asignar una influencia temporal a los antecedentes de cada pareja focal--rival. Los conteos y sus denominadores se agregarán de forma consistente: los aciertos de tiro se relacionarán con intentos y las tasas por posesión con sus posesiones elegibles. No se promediarán porcentajes sin considerar su exposición. Esta alternativa utilizará inicialmente un criterio de recencia compartido, sin diferencias por posición o variable. La ausencia de historial requerirá un tratamiento de referencia separado; asignar pesos no crea observaciones inexistentes.

\subsubsection{Tratamiento condicionado por posición}
\label{subsubsec:ponderacion-posicion}
Se considerará la posición para evitar que una referencia global oculte diferencias funcionales. El estudio de Gómez Ruano et al. aporta evidencia específica de diferencias estadísticas por puesto en la WNBA \cite{gomez2007posicioneswnba}. No demuestra que deba premiarse automáticamente el rebote de las interiores ni penalizarse su menor número de asistencias; fundamenta la comparación dentro de un contexto posicional.

Se incorporarán los macrorroles de la focal y de la rival como categorías diferenciadas. Se comparará una referencia común con referencias por Guard, Wing y Big; adicionalmente, podrá evaluarse un modelo con parámetros condicionados al macrorrol frente a modelos separados. No se adoptarán todas estas opciones simultáneamente. Esta configuración no introduce descuento temporal ni pesos de características adicionales. Los grupos con pocos registros recurrirán a referencias más generales. La correspondencia entre posiciones originales y macrorroles se documentará, sin asumir que las etiquetas originales sean idénticas a la clasificación propuesta.

\subsubsection{Ponderación por variable}
\label{subsubsec:ponderacion-variable}
Las características utilizadas para agrupar jugadoras pueden presentar escalas, redundancias y capacidades de diferenciación distintas. Huang et al. propusieron una extensión de k-means que estima pesos de variables durante el agrupamiento \cite{huang2005variableweighting}. Este antecedente sustenta estudiar pesos de características, pero no prueba que esos pesos identifiquen la aportación causal de una estadística al rendimiento deportivo.

Yamada y Fujii distinguieron perfiles de tiro y de rol ofensivo mediante conjuntos de características diferentes y posteriormente emplearon información de los grupos en modelos de eficiencia \cite{yamada2024offensivelineupanalysisbasketball}. En el presente proyecto se considerarán únicamente variables disponibles, sin presuponer acceso a tracking o a tipos de jugada anotados. Para caracterizar estilos individuales se construirá una vista por jugadora y fecha de corte, evitando que una actuación repetida en varias filas rivales se contabilice como evidencia individual independiente. El perfil resultante se incorporará después a las filas dirigidas de comparación.

Se contrastará el clustering con características estandarizadas y sin pesos adicionales frente a una selección acotada de pesos aprendidos. Esta operación no implica construir un indicador global de rendimiento ni ponderar la variable objetivo. El aprendizaje de relaciones en la regresión posterior seguirá siendo una decisión diferenciada.

\subsubsection{Combinación por temporada y posición}
\label{subsubsec:ponderacion-temporada-posicion}
Se evaluará si la información reciente conserva una pertinencia distinta según el macrorrol. La combinación reúne el antecedente de ponderación temporal y la evidencia de heterogeneidad posicional \cite{dixon1997scores,gomez2007posicioneswnba}; constituye una hipótesis del proyecto, no un procedimiento validado por esos trabajos para parejas rivales de la WNBA.

Como primera configuración se aplicará un descuento temporal común y se construirán referencias condicionadas al macrorrol. Solo si la cobertura lo permite se compararán descuentos diferentes para Guard, Wing y Big. Las columnas de temporada, fecha y macrorroles de ambas participantes permitirán distinguir estos contextos, manteniendo unidos los registros de cada partido. La fragmentación de la muestra y el incremento de parámetros deberán justificarse mediante resultados de validación.

\subsubsection{Combinación por temporada y variable}
\label{subsubsec:ponderacion-temporada-variable}
Se examinará si todas las métricas requieren el mismo horizonte histórico o si algunas se describen mejor con información más reciente. Los antecedentes de ponderación temporal y de pesos de características motivan esta comparación, pero no establecen qué indicadores son más persistentes en los datos del proyecto \cite{dixon1997scores,huang2005variableweighting}.

Se distinguirán dos experimentos. El primero construirá características con un descuento temporal común y posteriormente ponderará su contribución al clustering. El segundo permitirá descuentos temporales diferentes para un conjunto reducido de métricas, sin añadir simultáneamente pesos de características libres. En las filas focal--rival se conservarán los numeradores y denominadores necesarios para calcular cada antecedente. Esta separación permitirá identificar qué modificación aporta una mejora y limitar la cantidad de parámetros ajustables.

\subsubsection{Combinación por posición y variable}
\label{subsubsec:ponderacion-posicion-variable}
Se considerará que las características útiles para distinguir estilos pueden variar entre macrorroles. Las diferencias posicionales y la posibilidad de estimar pesos de características proporcionan antecedentes independientes para formular esta alternativa \cite{gomez2007posicioneswnba,huang2005variableweighting}. La asignación concreta de pesos deberá verificarse con los datos, sin imponer que una estadística sea importante únicamente por corresponder a una expectativa tradicional del puesto.

En la vista individual se comparará una representación común con otra que permita referencias o pesos de características por Guard, Wing y Big. Posteriormente, los perfiles de la focal y de la rival se incorporarán como categorías a cada comparación. Si el clustering se ajusta por separado en cada macrorrol, sus etiquetas se identificarán junto con el grupo de origen: un número de cluster no tendrá un significado universal ni un orden de calidad. Esta alternativa no incorporará descuento temporal adicional.

\subsubsection{Combinación por temporada, posición y variable}
\label{subsubsec:ponderacion-temporada-posicion-variable}
La configuración conjunta permitirá estudiar si la caracterización mejora al considerar simultáneamente recencia, macrorrol y contribución de las características. Los antecedentes anteriores sustentan sus componentes por separado; la utilidad de integrarlos constituye una propuesta experimental del presente trabajo.

Su evaluación quedará condicionada a que las alternativas simples demuestren una cobertura suficiente y un beneficio verificable. En el dataset se construirán antecedentes temporales por pareja, referencias posicionales y características individuales para el clustering. Los parámetros se compartirán entre grupos cuando resulte necesario y se limitarán las combinaciones ensayadas. No se ajustará un conjunto independiente de pesos para cada jugadora o pareja. Esta opción presenta la mayor carga de selección y podrá descartarse si sus mejoras son pequeñas o inestables.

\subsubsection{Comparación y criterios de selección}
\label{subsubsec:comparacion-ponderaciones}
La Tabla~\ref{tab:comparacion-ponderaciones-fase1} presenta una valoración metodológica preliminar de las alternativas. Las categorías de riesgo y complejidad son relativas al número de parámetros y a la fragmentación prevista de la muestra; no constituyen resultados experimentales ni propiedades universales de cada método.

\begingroup
\small
\setlength{\tabcolsep}{3pt}
\renewcommand{\arraystretch}{1.15}
\begin{longtable}{@{}p{2.4cm}p{2.85cm}p{1.95cm}p{2.05cm}p{3.15cm}@{}}
\caption{Comparación preliminar de las alternativas de ponderación y contextualización.}\label{tab:comparacion-ponderaciones-fase1}\\
\toprule
\textbf{Alternativa} & \textbf{Propósito principal} & \textbf{Comple\-jidad} & \textbf{Riesgo de sobreajuste} & \textbf{Limitación principal}\\
\midrule
\endfirsthead
\toprule
\textbf{Alternativa} & \textbf{Propósito principal} & \textbf{Comple\-jidad} & \textbf{Riesgo de sobreajuste} & \textbf{Limitación principal}\\
\midrule
\endhead
Temporada & Priorizar antecedentes recientes. & Baja & Bajo relativo & Puede perder información útil al descontar registros antiguos.\\
\midrule
Posición & Comparar dentro de contextos funcionales. & Media & Medio & Puede fragmentar la muestra y ocultar jugadoras versátiles.\\
\midrule
Variable & Diferenciar características del clustering. & Media & Medio & Los pesos pueden adaptarse al ruido o a características redundantes.\\
\midrule
Temporada y posición & Combinar recencia y referencias posicionales. & Media & Medio a alto & Los descuentos por grupo requieren suficiente historial.\\
\midrule
Temporada y variable & Diferenciar representación o recencia por métrica. & Media a alta & Medio a alto & Aumenta el número de configuraciones y dificulta atribuir mejoras.\\
\midrule
Posición y variable & Distinguir características dentro de cada macrorrol. & Alta & Alto relativo & Puede producir perfiles poco comparables entre grupos pequeños.\\
\midrule
Temporada, posición y variable & Integrar los tres ejes. & Alta & Alto relativo & Requiere mayor cobertura, control de parámetros e interpretación.\\
\bottomrule
\end{longtable}
\endgroup

Se conservará como referencia un procedimiento con escalamiento común y sin ponderaciones adicionales. Las alternativas se incorporarán de manera incremental y se evaluarán mediante estabilidad e interpretación de perfiles, error predictivo y cobertura de parejas. Una mejora del clustering no se considerará suficiente si no aporta utilidad a la predicción posterior. La elección se realizará con bloques cronológicos de validación y se comprobará una sola vez en el conjunto de prueba reservado.

Los parámetros de escalamiento, las referencias por posición, los pesos y el clustering se estimarán exclusivamente con información de entrenamiento anterior al momento de predicción \cite{sklearnCommonPitfallsPonderacion}. Esta restricción también se aplicará al enriquecimiento de la fase 2. El ajuste de hiperparámetros se desarrollará dentro de la etapa experimental prevista por la metodología en cascada y se documentará como parte de sus entregables.

\setcounter{secnumdepth}{\profundidadAntesPonderacion}
